\section{Matrices de validación}

Este anexo resume la lectura de validación usada para cerrar la memoria. Las tablas separan evidencia estadística, pruebas de consistencia y alcance de la inferencia. Las cifras se reportan como escala de evaluación y no como garantía de despliegue industrial.

\begin{table}[H]
\centering
\scriptsize
\begin{tabularx}{\textwidth}{@{}p{0.10\textwidth}p{0.26\textwidth}p{0.22\textwidth}X@{}}
\toprule
Estudio & Magnitud validada & Escala & Decisión de lectura \\
\midrule
SP1 & Reclutamiento y satisfacción de demanda. & $27\,300$ evaluaciones. & Evidencia favorable para coaliciones distribuidas; la superioridad depende de métrica y coste. \\
SP2 & Capacidad efectiva heterogénea. & $20\,280$ comprobaciones. & La capacidad marginal mejora la lectura frente a cardinalidad plana. \\
SP3 & Residual de fuerza y torque. & $20\,040$ corridas y $38\,076$ comprobaciones. & Se refuta la equivalencia entre cobertura escalar y factibilidad física. \\
SP4 & Movimiento, llegada y seguridad. & $20\,070$ ejecuciones y $21\,408$ comprobaciones. & El resultado debe leerse como compromiso Pareto, no como ranking absoluto. \\
SP5 & Transporte de carga extendida. & $20\,040$ ejecuciones/comprobaciones. & La pose final de la carga y los despejes cambian la evaluación respecto a llegada individual. \\
SP6 & Recuperación ante eventos. & $20\,000$ ejecuciones/comprobaciones. & La robustez exige medir reparación, tiempo, carga perdida y seguridad posterior. \\
SP7 & Comunicación y sensado degradado. & $20\,176$ evaluaciones/comprobaciones. & La conectividad temporal es más informativa que la conectividad instantánea. \\
SP8 & Escala de flota e intractabilidad. & $20\,000$ evaluaciones hasta $50\,000$ AMR. & La escala favorece métodos reparables y distribuidos; el modelo sigue siendo mesoscópico. \\
\bottomrule
\end{tabularx}
\caption{Matriz de validación por estudio experimental.}
\label{tab:matriz-validacion-sp}
\end{table}

\subsection*{C.1 Mapa extendido de evidencia}

La lectura canónica de SP1--SP8 se apoya en barridos complementarios. La tabla siguiente
separa qué añade cada capa de evaluación para que las conclusiones no dependan de una
sola parametrización ni de un único tipo de escenario.

\begin{table}[H]
\centering
\scriptsize
\begin{tabularx}{\textwidth}{@{}p{0.08\textwidth}p{0.28\textwidth}p{0.24\textwidth}X@{}}
\toprule
Estudio & Capas evaluadas & Profundidad & Función dentro de la validación \\
\midrule
SP1 & Reclutamiento principal, comparación con referencia central y control visual de trayectorias. & $27\,300$ evaluaciones, $65$ ordenaciones y $51$ inspecciones cinemáticas. & Fija la base de comparación: cerrar demanda con coste razonable antes de introducir capacidad física. \\
SP2 & Capacidad efectiva y ablación de utilidad marginal. & $20\,280$ comprobaciones principales, más $6$ familias de ablación. & Aísla el efecto de ponderar aportes heterogéneos frente a contar robots como unidades equivalentes. \\
SP3 & Barrido de \emph{wrench}, versión de alta exigencia y batería dinámica de pose. & $20\,040$ corridas, $38\,076$ comprobaciones y validación con modelo Euler--Lagrange. & Demuestra que fuerza y torque deben evaluarse como magnitudes vectoriales, no como capacidad escalar acumulada. \\
SP4 & Ley explícita de movimiento, comparación de navegación y variante de alta potencia. & $72$ ejecuciones de ley cerrada, $300$ casos pareados y $20\,070$ ejecuciones exigentes. & Separa llegada, seguridad, energía y despeje; evita declarar superioridad por una sola métrica de movimiento. \\
SP5 & Transporte cooperativo, variante de alta potencia y ley explícita aplicada a carga extendida. & $20\,040$ ejecuciones principales, $334$ semillas y validación específica de control. & Verifica que la pose de la carga, la formación y la envolvente geométrica cambian la lectura frente a llegada individual. \\
SP6 & Recuperación explícita, robustez comparativa y variante de alta potencia. & $20\,000$ ejecuciones principales, $250$ semillas y $8$ familias de escenario. & Mide reparación tras fallo, batería baja, retardo, carga inviable y seguridad posterior al evento. \\
SP7 & Comunicación degradada con exigencia estadística. & $20\,176$ evaluaciones, $97$ semillas y perfiles de radio, pérdida, retardo, jitter y sensado. & Comprueba que el transporte depende de conectividad temporal, descubrimiento local y tolerancia a pérdida de mensajes. \\
SP8 & Escalera de flota, variante de alta potencia y escenario de almacén. & $20\,000$ evaluaciones, $20$ tamaños de flota y escala hasta $50\,000$ AMR. & Delimita cuándo una referencia central pierde practicidad y cuándo una política distribuida conserva relación útil entre calidad y recurso. \\
\bottomrule
\end{tabularx}
\caption{Mapa extendido de la evidencia experimental SP1--SP8.}
\label{tab:mapa-extendido-sp}
\end{table}

\begin{table}[H]
\centering
\scriptsize
\begin{tabularx}{\textwidth}{@{}p{0.18\textwidth}p{0.30\textwidth}X@{}}
\toprule
Tipo de evidencia & Qué permite afirmar & Qué no permite afirmar \\
\midrule
Prueba algebraica & Coherencia interna del modelo, existencia de contraejemplos y necesidad de variables físicas. & Rendimiento estadístico bajo escenarios variados. \\
Monte Carlo pareado & Diferencias agregadas entre métodos bajo mundos compartidos y semillas controladas. & Garantía universal fuera del dominio simulado. \\
Inspección visual & Plausibilidad de trayectorias, formación, recuperación y movimiento de carga. & Validación física en hardware o contacto realista completo. \\
Referencia centralizada & Cota superior o punto de comparación para medir coste de descentralización. & Solución operativa a gran escala cuando el coste se vuelve prohibitivo. \\
Método aprendido & Comparación adaptativa frente a reglas analíticas y referencias clásicas. & Garantía estructural, interpretabilidad completa o dominancia fuera de distribución. \\
\bottomrule
\end{tabularx}
\caption{Alcance de cada tipo de evidencia usado en la memoria.}
\label{tab:alcance-evidencia}
\end{table}

\begin{table}[H]
\centering
\scriptsize
\begin{tabularx}{\textwidth}{@{}p{0.18\textwidth}X X@{}}
\toprule
Riesgo de sobreafirmación & Formulación defendible & Formulación que se evita \\
\midrule
Capacidad \emph{wrench} & Certificado planar cuasiestático para decidir si una coalición asignada puede generar fuerza y torque. & Prueba completa de manipulación física real con fricción y contacto 3D. \\
Escalabilidad & Comparación mesoscópica de coste, completitud y factibilidad hasta escalas de almacén. & No afirmar validación industrial directa en una instalación real. \\
Vídeos & Evidencia cualitativa de inspección asociada a métricas. & Prueba suficiente por sí sola. \\
Métodos aprendidos & Comparadores entrenados bajo el simulador y evaluados con separación de fases. & Estado del arte universal o solución dominante. \\
Seguridad & Medición de colisiones, despejes y restricciones en el dominio simulado. & Certificación de seguridad funcional. \\
\bottomrule
\end{tabularx}
\caption{Reglas de redacción para mantener afirmaciones defendibles.}
\label{tab:riesgos-sobreafirmacion}
\end{table}

\IfFileExists{../generated/stats_master_table.tex}{% Auto-generated. Do not edit by hand.
\begin{table}[htbp]
\centering
\scriptsize
\resizebox{\textwidth}{!}{%
\begin{tabular}{lllllll}
\toprule
SP & Hipotesis & Metrica & n & p-Holm & Efecto & Veredicto \\
\midrule
SP1 & H01\_methods\_differ\_theoretical\_gap & optimality\_gap\_vs\_oracle & 2100 & 0 & 0.6063 & rechaza H0 \\
SP1 & H02\_mappo\_lower\_gap\_than\_imitation & optimality\_gap\_vs\_oracle & 2100 & 0 & -0.8951 & rechaza H0 \\
SP1 & H03\_mappo\_lower\_gap\_than\_hungarian\_classic & optimality\_gap\_vs\_oracle & 2100 & 1.56e-182 & -0.6904 & rechaza H0 \\
SP1 & H04\_smith\_lower\_runtime\_than\_mappo & runtime\_ms & 2100 & 0 & -1.092 & rechaza H0 \\
SP1 & H05\_mappo\_higher\_coalition\_success\_than\_smith & coalition\_success\_rate & 2100 & 0 & 1.482 & rechaza H0 \\
SP2 & H01\_methods\_differ\_capacity\_gap & optimality\_gap\_vs\_oracle & 1560 & 0 & 0.71 & rechaza H0 \\
SP2 & H01b\_methods\_differ\_capacity\_ceiling\_gap & capacity\_gap\_vs\_capacity\_oracle & 1560 & 5.68e-10 & 0.008093 & rechaza H0 \\
SP2 & H02\_primal\_dual\_higher\_capacity\_success\_than\_greedy & capacity\_success\_rate & 1560 & 1 & -0.7076 & no rechaza H0 \\
SP2 & H03\_neural\_lower\_gap\_than\_linear\_imitation & optimality\_gap\_vs\_oracle & 1560 & 1.17e-26 & -0.2328 & rechaza H0 \\
SP2 & H04\_local\_primal\_dual\_lower\_messages\_than\_oracle & communication\_messages & 1560 & 6.85e-18 & -0.362 & rechaza H0 \\
SP2 & H05\_smith\_lower\_runtime\_than\_neural & runtime\_ms & 1560 & 8.71e-256 & -1.589 & rechaza H0 \\
SP3 & H3\_HP\_1\_scalar\_false\_positives\_rotational\_loads & false\_positive\_rate & 2004 & 8.47e-198 & 0.9998 & rechaza H0 \\
SP3 & H3\_HP\_2\_smith\_wrench\_lower\_residual\_than\_capacity & wrench\_residual\_feasible\_available & 2004 & 1 & -0.0316 & no rechaza H0 \\
SP3 & H3\_HP\_3\_capacity\_greedy\_more\_infeasible\_assignments\_than\_support\_dual & fp\_given\_assigned & 2004 & 1 & 0.0316 & no rechaza H0 \\
SP3 & H3\_HP\_4\_support\_dual\_lower\_gap\_than\_cbba & optimality\_gap\_vs\_wrench\_oracle & 2004 & 1 & 0.02293 & no rechaza H0 \\
SP3 & H3\_HP\_5\_methods\_differ\_wrench\_gap & optimality\_gap\_vs\_wrench\_oracle & 2004 & 0 & 0.1175 & rechaza H0 \\
SP4 & H4\_HP\_1\_reference\_reduces\_collision\_vs\_direct & collision\_rate & 1338 & 2.52e-183 & -0.8525 & rechaza H0 \\
SP4 & H4\_HP\_2\_cbf\_reduces\_collision\_vs\_direct & collision\_rate & 1338 & 8.88e-183 & -0.8471 & rechaza H0 \\
SP4 & H4\_HP\_3\_tensor\_reduces\_collision\_vs\_direct & collision\_rate & 1338 & 1.41e-183 & -0.843 & rechaza H0 \\
SP4 & H4\_HP\_4\_energy\_smith\_reduces\_energy\_vs\_cbf & energy\_proxy\_wh & 1338 & 1 & -0.06049 & no rechaza H0 \\
SP4 & H4\_HP\_5\_methods\_differ\_gap\_vs\_reference & performance\_gap\_vs\_reference & 1338 & 0 & 0.3687 & rechaza H0 \\
SP5 & H5\_1\_reference\_reduces\_gap\_vs\_classic\_apf & performance\_gap\_vs\_reference & 2004 & 0 & -31.63 & rechaza H0 \\
SP5 & H5\_2\_ours\_hamiltonian\_reduces\_formation\_break\_vs\_sota\_vo\_cargo & formation\_broken\_rate & 2004 & 0.2153 & -0.06775 & no rechaza H0 \\
SP5 & H5\_3\_ours\_tensor\_improves\_target\_reached\_vs\_classic\_push & target\_reached & 2004 & 0.5356 & -0.0316 & no rechaza H0 \\
SP5 & H5\_4\_methods\_differ\_transport\_score & score\_value & 2004 & 0 & 0.7322 & rechaza H0 \\
SP6 & H6.1\_ours\_reduces\_lost\_load\_rate\_vs\_classic\_greedy & lost\_load\_rate & 2000 & 2.55e-05 & -0.1359 & rechaza H0 \\
SP6 & H6.2\_ours\_improves\_completion\_vs\_smith\_qr & task\_completion\_rate & 2000 & 0.1962 & 0.03077 & no rechaza H0 \\
SP6 & H6.3\_ours\_reduces\_reference\_gap\_vs\_cbba & performance\_gap\_vs\_reference & 2000 & 2.06e-92 & -0.3084 & rechaza H0 \\
SP6 & H6.4\_recovery\_family\_differs\_on\_reference\_gap & performance\_gap\_vs\_reference & 2000 & 7.19e-160 & 0.07529 & rechaza H0 \\
SP7 & H7.1\_radius\_improves\_coalition\_connectivity & coalition\_connected\_time\_ratio & 20176 & 0 & 0.8934 & rechaza H0 \\
SP7 & H7.2\_ours\_connectivity\_beats\_classic\_under\_harsh\_profiles & transport\_network\_score & 485 & 1.36e-80 & 0.8804 & rechaza H0 \\
SP7 & H7.3\_methods\_differ\_under\_intermittent\_communication & transport\_network\_score &  & 1 & nan & no rechaza H0 \\
SP8 & H8.1\_centralized\_oracle\_timeout\_increases\_with\_scale & timeout\_rate & 2000 & 4.93e-287 & 0.6937 & rechaza H0 \\
SP8 & H8.2\_ours\_hierarchical\_higher\_completion\_than\_classic\_local & task\_completion\_rate & 2000 & 4.18e-300 & 0.7791 & rechaza H0 \\
SP8 & H8.3\_ours\_hierarchical\_higher\_wrench\_feasibility\_than\_hungarian & wrench\_feasible\_rate & 2000 & 6.45e-299 & 1.167 & rechaza H0 \\
\bottomrule
\end{tabular}
}%
\caption{Tabla estadistica consolidada generada desde hypothesis\_results.csv canonicos.}
\label{tab:stats-master-generated}
\end{table}
}{}
